%%
%% spring-lecture13.tex
%% 
%% Made by Alex Nelson <pqnelson@gmail.com>
%% Login   <alex@lisp>
%% 
%% Started on  2026-04-30T11:31:04-0700
%% Last update 2026-04-30T11:31:04-0700
%% 

\lecture{}

\begin{node}
We will prove the Krull--Schmidt theorem (the most important theorem
in this part of the course) and then pivot to representation theory of
algebras. 
\end{node}

\begin{lemma}\label{lemma:5.7}
For any short exact sequence of $A$-modules
\begin{equation}
0\to N\xrightarrow{\phi}M\xrightarrow{\psi}P\to0,
\end{equation}
the following are equivalent:
\begin{enumerate}
\item There exists a $\chi\colon P\to M$ such that $\psi\circ\chi=\id_{P}$
\item There exists a $\theta\colon M\to N$ such that $\theta\circ\phi=\id_{N}$
\end{enumerate}
In these cases,
$M=\Im(\chi)\oplus\ker(\psi)=\Im(\phi)\oplus\ker(\theta)$. We call such
an exact sequence \define{Split}, $\psi$ is a ``split surjection'',
$\phi$ is a ``split injection''.
\end{lemma}

\begin{proof}
$(2)\implies(1)$ is homework.

$(1)\implies(2)$ For any $u\in M$,
\begin{equation}
u=\chi(\psi(u))+(u-\chi(\psi(u)))
\end{equation}
and
\begin{equation}
\psi(u-\chi(\psi(u)))=\psi(u)-\psi(\chi(\psi(u)))=0,
\end{equation}
so $M=\Im(\chi)+\ker(\psi)$.

We want to show $\Im(\chi)\cap\ker(\psi)=0$ (to make it a direct
sum). Since $\psi\circ\chi=\id$,
\begin{equation}
\ker(\psi)\cap\Im(\chi)=0
\end{equation}
[Pf: otherwise there's a nonzero $0\neq x\in P$ such that $x=\psi(\chi(x))=0$
  which is a contradiction].
Now because
\begin{equation}\label{eq:prop5.7:pf:star}
\Im(\phi)=\ker(\psi)\quad\mbox{and $\phi$ is injective},
\end{equation}
the mapping
\begin{equation}
\theta\colon M\to N
\end{equation}
such that $\theta(u)=\phi^{-1}(u-\chi(\psi(u)))$ is a morphism an
satisfies for every $v\in N$,
\begin{subequations}
  \begin{align}
\theta(\phi(v)) &= \phi^{-1}(\phi(v)-\chi(\psi(\phi(v))))\\
&=v-\phi^{-1}(\chi(\psi(\phi(v))))+\ker(\phi)\\
&=v-\phi^{-1}(\chi(\psi(\phi(v))))+0\\
&=v\quad\mbox{by Eq~\eqref{eq:prop5.7:pf:star}}
  \end{align}
\end{subequations}
so $\theta\circ\phi=\id_{N}$.
\end{proof}

\begin{lemma}\label{lemma:5.8}
Let $M=M_{1}\oplus M_{2}=N_{1}\oplus N_{2}$ be decompositions of the
$A$-module $M$. Assume there exists $\phi\in\Aut_{A}(M)$ such that we
can write it in matrix form
\begin{equation}
\phi=[\phi_{i,j}]=\begin{bmatrix}\phi_{1,1} & \phi_{1,2}\\\phi_{2,1} & \phi_{2,2}
\end{bmatrix}
\end{equation}
If $\phi_{1,1}\colon M_{1}\to N_{1}$ is an isomorphism, then
$M_{2}\iso N_{2}$.
\end{lemma}

\begin{proof}
For
\begin{equation}
\phi'=\begin{bmatrix}\id_{N_{1}} & 0\\
-\phi_{21}\circ\phi_{11}^{-1} & \id_{N_{2}}
\end{bmatrix}
\end{equation}
and
\begin{equation}
\phi''=\begin{bmatrix}\id_{M_{1}} & -\phi_{11}^{-1}\circ\phi_{12}\\
0 & \id_{M_{2}}
\end{bmatrix}
\end{equation}
these are automorphisms of $M$. Then $\phi'\circ\phi\circ\phi''$ is an
automorphism equal to
\begin{equation}
\phi'\circ\phi\circ\phi''=\begin{bmatrix}\phi_{11} & 0\\
0 & \psi\end{bmatrix}
\end{equation}
where
\begin{equation}
\psi=\phi_{22}-\phi_{21}\circ\phi_{11}^{-1}\circ\phi_{12}\colon
M_{2}\to N_{2}
\end{equation}
is an isomorphism.
\end{proof}

\begin{proposition}[Krull--Schmidt]\label{prop:5.9}
Let $M=\bigoplus^{r}_{i=1}M_{i}$ and $N=\bigoplus^{s}_{j=1}N_{j}$ be
such that $\End_{A}M_{i}$ and $\End_{A}N_{j}$ are local algebras. If
$M\iso N$, then $r=s$ and there exists a permutation $\sigma\in S_{r}$
such that $M_{i}=N_{\sigma(i)}$.
\end{proposition}

\begin{corollary}\label{cor:5.10}
\begin{enumerate}
\item If $M$ is Artininan and Noetherian, then it has the unique
  decomposition into indecomposible modules.
\item If $A$ is right Artinian, then every finitely-generated
  $A$-module uniquely decomposes into a finite direct sum of
  indecomposable modules.
\end{enumerate}
\end{corollary}

\begin{proof}
\begin{enumerate}
\item By Proposition~\ref{prop:5.1}, Corollary~\ref{cor:5.6}, and Proposition~\ref{prop:5.9}.
\item By Corollary~\ref{cor:4.16} and Corollary~\ref{cor:5.10}~(1).\qedhere
\end{enumerate}
\end{proof}

\subsection{Representations of Algebras}

\begin{definition}
Let $\FF$ be a field, let $A$ be an $\FF$-algebra. A \define{(Matrix) Representation}
of $A$ is an $\FF$-algebra morphism $\theta\colon A\to M_{n}(\FF)$.

The \define{Degree} of the representation $\theta\colon A\to M_{n}(\FF)$
is $\deg(\theta)=n$.
\end{definition}

\begin{node}
We will just speak of ``representations'' of $A$, understanding that
we mean ``matrix representations'' of $A$. In particular, they are
\define{Finite-Dimensional} representations because
$\dim_{\FF}(M_{n}(\FF))=n^{2}<\infty$. 
\end{node}

\begin{definition}
Let $\theta\colon A\to M_{n}(\FF)$ be a representation of $A$. We say
the representation $\theta$ is \define{Faithful} if $\ker(\theta)=0$
(so $\dim_{\FF}A\leq n^{2}$). In particular, if $\dim_{\FF}A=\infty$,
there are no faithful matrix representations of $A$.
\end{definition}

\begin{definition}
Let $\theta$, $\phi$ be representations of $A$ with $m=\deg(\phi)$ and
$n=\deg(\theta)$. Let $\alpha\in M_{n,m}(\FF)$ be an $n\times m$
matrix. We say $\alpha$ is an \define{Intertwiner} if for all $x\in A$
we have $\theta(x)\alpha=\alpha\phi(x)$. In this case, we write
\begin{equation}
\alpha\colon\theta\to\phi.
\end{equation}
We should intuitively think of an intertwiner as a triple
$(\theta,\alpha,\psi)$ since the same matrix may intertwine different
representations. 
\end{definition}

\begin{fact}
The $\FF$-matrix representations of $A$ and their intertwiners form a
category which we will denote by $\MatRep_{\FF}(A)$. This is a
subcategory of functors from the delooping $\cat{B}A$ of $A$ to the
category $\cat{Mat}_{\FF}$ whose objects are $M_{n}(\FF)$ matrix
algebras, and morphisms are $\FF$-linear transformations
(multiplication by an $m\times n$ matrix).

More generally, we can replace $\cat{Mat}_{\FF}$ with any other
suitable category, and we get other types of reprensetations for
an associative algebra.
\end{fact}

\begin{definition}
Let $\theta$, $\phi$ be $\FF$-matrix representations of $A$. We say
$\theta$ and $\phi$ are \define{Equivalent} if they are isomorphic in
the category theoretic sense, i.e., there exists intertwiners
$\alpha\colon\theta\to\phi$ and $\beta\colon\phi\to\theta$ such that
$\alpha\circ\beta=\id_{\phi}$ and $\beta\circ\alpha=\id_{\theta}$.

In this case, we write $\theta\iso\phi$ to indicate they are equivalent.
\end{definition}

\begin{xca}
Prove $\theta\iso\phi$ iff $\deg\theta=\deg\phi$ and there exists a
nonsingular intertwiner $\alpha\colon\theta\to\phi$.
\end{xca}

\begin{definition}
Let $\theta$ be an $\FF$-representation of $A$, write
$M_{\theta}=\bigoplus n\FF$ (where $n=\deg(\theta)$) equipped with the 
$A$-action
\begin{equation}
[a_{1},\dots,a_{n}]x:=[a_{1},\dots,a_{n}]\theta(x)
\end{equation}
where $[a_{1},\dots,a_{n}]\in M_{\theta}$ is a row vector, $x\in A$.
This defines an $A$-module such that
$\dim_{\FF}M_{\theta}=\deg(\theta)$ and $\Annihilator(M_{\theta})=\ker(\theta)$.

The correspondence $M\colon\theta\mapsto M_{\theta}$ sending
$\FF$-representations [of a fixed algebra $A$] to $A$-modules, and
$\alpha\mapsto\mu_{\alpha}$ where $\alpha\colon\theta\to\psi$ is an
intertwiner and the linear operator $\mu_{\alpha}\colon M_{\theta}\to M_{\psi}$ sends
$[a_{1},\dots,a_{n}]\mapsto[a_{1},\dots,a_{n}]\alpha$ --- this forms a
functor $\MatRep_{\FF}(A)\to\Mod(A)$ from the category of matrix
representations of $A$ to the category of $A$-modules.
\end{definition}

